\section{Appendix}

This appendix collects supporting analyses that are useful for interpretation
but too detailed for the six-page body.

\section{Additional Cross-Lingual Retrieval Diagnostics}

These diagnostics expand the paper's main cross-language retrieval result. They
show that the same-language advantage is not a single-score artifact: it appears
as model-level home advantage, directed language-pair asymmetry, delayed recovery
of multilingual publication variants, and weaker embedding separability for
cross-language gold evidence.

\begin{figure}[t]
  \centering
  \includegraphics[width=\linewidth]{cp_fig02_home_advantage.png}
  \caption{Same-language home advantage by model. The gap between
  same-language and cross-language retrieval is the main reason aggregate recall
  can overstate deployment readiness.}
  \label{fig:app-home-advantage}
\end{figure}

\begin{figure}[t]
  \centering
  \includegraphics[width=\linewidth]{cp_fig03_directional_clir_matrix.png}
  \caption{Directional query-language $\rightarrow$ document-language retrieval.
  Cross-language retrieval is not symmetric: a model's behavior depends on both
  the query language and the target document language.}
  \label{fig:app-directional-clir}
\end{figure}

\begin{figure}[t]
  \centering
  \includegraphics[width=\linewidth]{cp_fig06_07_mate.png}
  \caption{Foreign-twin retrieval diagnostics. The figure reports whether
  multilingual publication variants are found and how deeply a model must search
  before reaching the first cross-language twin.}
  \label{fig:app-mate-retrieval}
\end{figure}

\begin{figure}[t]
  \centering
  \includegraphics[width=\linewidth]{cp_fig11_separability.png}
  \caption{Embedding separability for same-language and cross-language gold
  documents. Lower cross-language separability indicates that the retrieval
  failure is partly representation-level, not only a top-$k$ ranking artifact.}
  \label{fig:app-separability}
\end{figure}

\section{Additional Alias-Graph Diagnostics}

The alias-graph analyses isolate the chemistry side of the retrieval problem.
They test whether multilingual queries for the same concept produce consistent
rankings and whether chemically plausible neighbors become systematic
distractors. Together, these figures explain why relevance in chemistry search
cannot be reduced to topical similarity alone.

\begin{figure}[t]
  \centering
  \includegraphics[width=\linewidth]{ag_fig1_cross_lingual_rbo.png}
  \caption{Cross-lingual ranking consistency for the same chemical concept asked
  in multiple languages. Low overlap means language variants of the same query
  can return substantially different ranked lists.}
  \label{fig:app-alias-rbo}
\end{figure}

\begin{figure}[t]
  \centering
  \includegraphics[width=\linewidth]{ag_fig12_joint_failure_modes.png}
  \caption{Joint language and chemistry failure modes among confused cases. This
  analysis links the paper's two central risks: same-language preference and
  chemically plausible but wrong neighbors.}
  \label{fig:app-joint-failures}
\end{figure}

\begin{figure}[t]
  \centering
  \includegraphics[width=\linewidth]{ag_fig6_question_type_effect.png}
  \caption{Question-type effects in the alias-graph benchmark. Some query types
  are substantially more vulnerable to chemistry-confusability than others.}
  \label{fig:app-question-type}
\end{figure}

\begin{figure}[t]
  \centering
  \includegraphics[width=\linewidth]{ag_fig8_confusion_is_separability.png}
  \caption{Confusion as separability collapse. Confused examples have weaker
  embedding separation between gold and non-gold documents.}
  \label{fig:app-confusion-separability}
\end{figure}

\begin{figure}[t]
  \centering
  \includegraphics[width=\linewidth]{ag_fig5_universal_attractors.png}
  \caption{Universal attractors in the alias graph: chemically plausible
  look-alikes that repeatedly steal top rank across models or languages.}
  \label{fig:app-universal-attractors}
\end{figure}

\section{Deployment-Oriented Supporting Analyses}

These figures support the deployment interpretation in the main paper. They are
secondary to the main cost-vs-capability decision, but they show two practical
constraints: query translation can change retrieval behavior, and re-rankers can
only recover cross-language evidence that appears within a feasible candidate
pool.

\begin{figure}[t]
  \centering
  \includegraphics[width=\linewidth]{cp_fig05_mt_penalty.png}
  \caption{Human-written versus machine-translated queries. This supports the
  paper's practical separation between generating query variants and relying on
  human-translated source evidence.}
  \label{fig:app-mt-penalty}
\end{figure}

\begin{figure}[t]
  \centering
  \includegraphics[width=\linewidth]{cp_fig19_rrc_budget.png}
  \caption{Re-ranker recoverability curves. These curves show how much
  cross-language evidence is recoverable at different candidate-pool depths and
  why some failures must be fixed before re-ranking.}
  \label{fig:app-rrc-budget}
\end{figure}
